% test-011.tex -- comandos \spnZ, \spmQ, \sprazon y \spProp
%
% Compilar:  lualatex-dev test-011.tex
% Validar:   show-pdf-tags --xml test-011.pdf | rnv-wrapp
%            verapdf --flavour ua2 --format text test-011.pdf
%
% Cada caso es un parrafo numerado, para poder referirse a el al escuchar
% el PDF. Los comentarios "doc:" indican lo que documenta el manual; en el
% resto, la lectura se revisa escuchando. Todos trabajan en modo
% matematico. \spnZ y el entero de \spmQ solo reciben enteros; en \spmQ el
% numerador tambien debe ser entero (incluso con use-spnZ=false) y el
% denominador distinto de cero. \sprazon y \spProp reciben cualquier
% expresion matematica. Las entradas invalidas (que detienen la
% compilacion) no van aqui. Las lecturas de \sprazon y \spProp se deducen
% de sus intents y hay que confirmarlas escuchando.
\DocumentMetadata
  {
    lang=es-CL,
    pdfversion=2.0,
    pdfstandard=ua-2,
    tagging=on,
    tagging-setup={math/setup={mathml-SE}},
  }
\documentclass{article}
\usepackage[spanish]{babel}
\usepackage{unicode-math}
\usepackage{spintent}
\usepackage{hyperref}
\hypersetup
  {
    pdfdisplaydoctitle = true,
    pdftitle           = {test-011: spnZ, spmQ, sprazon y spProp},
  }
\pagestyle{empty}
\begin{document}

\section*{A. Comando spnZ: argumentos}

% doc: «23»
1. Entero positivo: $\spnZ{23}$.

% doc: «menos 23»
2. Con signo negativo: $\spnZ{-23}$.

% doc: «más 23»
3. Con signo positivo explícito: $\spnZ{+23}$.

% doc: «0»
4. Cero: $\spnZ{0}$.

% doc: «100»
5. Cien: $\spnZ{100}$.

% doc: «1000000»
6. Un millón: $\spnZ{1000000}$.

% doc: «1234567»
7. Siete cifras: $\spnZ{1234567}$.

% doc: «menos 1000»
8. Negativo grande: $\spnZ{-1000}$.

\section*{B. Comando spnZ: llaves}

% doc: «1000000»
9. Con cifras-sep false: $\spnZ[cifras-sep=false]{1000000}$.

% doc: «100000»
10. Con cifras-sep false y espacios del usuario: $\spnZ[cifras-sep=false]{10\,00\,00}$.

% doc: «23 negativo»
11. Con read-sign clear y signo negativo: $\spnZ[read-sign=clear]{-23}$.

% doc: «23 positivo»
12. Con read-sign clear y signo positivo: $\spnZ[read-sign=clear]{+23}$.

% doc: «menos 23»
13. Con wrap-parens false, el valor por defecto: $\spnZ[wrap-parens=false]{-23}$.

% doc: «se abren paréntesis menos 23 se cierran paréntesis»
14. Con wrap-parens true: $\spnZ[wrap-parens=true]{-23}$.

% doc: «se abren paréntesis menos 23 se cierran paréntesis»
15. Con wrap-parens true y read-parens true: $\spnZ[wrap-parens=true,read-parens=true]{-23}$.

% doc: «menos 23»
16. Con wrap-parens true y read-parens false: $\spnZ[wrap-parens=true,read-parens=false]{-23}$.

% doc: «menos 23»
17. Con wrap-parens false y read-parens false: $\spnZ[wrap-parens=false,read-parens=false]{-23}$.

% doc: «1234567 negativo»
18. Con todas las llaves: $\spnZ[cifras-sep=false,read-sign=clear,wrap-parens=true,read-parens=false]{-1234567}$.

\section*{C. Comando spnZ: en fórmulas}

% doc: «menos 3 más 5»
19. Suma de enteros: $\spnZ{-3} + \spnZ{5}$.

% doc: «se abren paréntesis menos 3 se cierran paréntesis por 4»
20. Con paréntesis y producto: $\spnZ[wrap-parens=true]{-3} \spmsym \spnZ{4}$.

% doc: «12 dividido 4 es igual a 3»
21. División con igualdad: $\spnZ{12} \spdsym \spnZ{4} = \spnZ{3}$.

% doc: «se abren paréntesis menos 3 se cierran paréntesis por se abren paréntesis menos 4 se cierran paréntesis es igual a 12»
22. En fórmula destacada:
\[ \spnZ[wrap-parens=true]{-3} \spmsym \spnZ[wrap-parens=true]{-4} = \spnZ{12} \].

\section*{D. Comando spmQ: ejemplos del manual y llave join-word}

% doc: «2 y la fracción con numerador 5 y denominador 7»
23. Por defecto: $\spmQ{2}{5}{7}$.

% doc: «1 entero la fracción con numerador 5 y denominador 7»
24. Con join-word entero: $\spmQ[join-word=entero]{1}{5}{7}$.

% doc: «2 enteros la fracción con numerador 5 y denominador 7»
25. Con join-word enteros: $\spmQ[join-word=enteros]{2}{5}{7}$.

% doc: «2 y la fracción con numerador 5 y denominador 7»
26. Con join-word y: $\spmQ[join-word=y]{2}{5}{7}$.

\section*{E. Comando spmQ: denominadores}

% doc: «1 y la fracción con numerador 1 y denominador 2»
27. Con denominador 2: $\spmQ{1}{1}{2}$.

% doc: «1 y la fracción con numerador 1 y denominador 3»
28. Con denominador 3: $\spmQ{1}{1}{3}$.

% doc: «2 y la fracción con numerador 3 y denominador 4»
29. Con denominador 4: $\spmQ{2}{3}{4}$.

% doc: «3 y la fracción con numerador 2 y denominador 5»
30. Con denominador 5: $\spmQ{3}{2}{5}$.

% doc: «1 y la fracción con numerador 1 y denominador 6»
31. Con denominador 6: $\spmQ{1}{1}{6}$.

% doc: «4 y la fracción con numerador 6 y denominador 7»
32. Con denominador 7: $\spmQ{4}{6}{7}$.

% doc: «5 y la fracción con numerador 3 y denominador 8»
33. Con denominador 8: $\spmQ{5}{3}{8}$.

% doc: «6 y la fracción con numerador 1 y denominador 9»
34. Con denominador 9: $\spmQ{6}{1}{9}$.

% doc: «7 y la fracción con numerador 3 y denominador 10»
35. Con denominador 10: $\spmQ{7}{3}{10}$.

% doc: «8 y la fracción con numerador 1 y denominador 11»
36. Con denominador 11: $\spmQ{8}{1}{11}$.

% doc: «9 y la fracción con numerador 5 y denominador 12»
37. Con denominador 12: $\spmQ{9}{5}{12}$.

% doc: «3 y la fracción con numerador 1 y denominador 100»
38. Con denominador 100: $\spmQ{3}{1}{100}$.

\section*{F. Comando spmQ: enteros y llaves}

% doc: «menos 2 y la fracción con numerador 5 y denominador 7»
39. Entero negativo: $\spmQ{-2}{5}{7}$.

% doc: «más 2 y la fracción con numerador 5 y denominador 7»
40. Entero con signo positivo explícito: $\spmQ{+2}{5}{7}$.

% doc: «100 y la fracción con numerador 1 y denominador 2»
41. Entero grande: $\spmQ{100}{1}{2}$.

% doc: «1000 y la fracción con numerador 1 y denominador 2»
42. Entero de cuatro cifras: $\spmQ{1000}{1}{2}$.

% doc: «2 y la fracción con numerador 5 y denominador 7»
43. Con use-spnZ true, el valor por defecto: $\spmQ[use-spnZ=true]{2}{5}{7}$.

% doc: «2 y 5 séptimos»
44. Con use-spnZ false: $\spmQ[use-spnZ=false]{2}{5}{7}$.

% doc: «2 y la fracción con numerador 5 y denominador 7»
45. Con use-dfrac false, el valor por defecto: $\spmQ[use-dfrac=false]{2}{5}{7}$.

% doc: «2 y la fracción con numerador 5 y denominador 7»
46. Con use-dfrac true: $\spmQ[use-dfrac=true]{2}{5}{7}$.

% doc: «se abren paréntesis 2 y la fracción con numerador 5 y denominador 7 se cierran paréntesis»
47. Con wrap-parens true: $\spmQ[wrap-parens=true]{2}{5}{7}$.

% doc: «2 y la fracción con numerador 5 y denominador 7»
48. Con wrap-parens true y read-parens false: $\spmQ[wrap-parens=true,read-parens=false]{2}{5}{7}$.

% doc: «2 enteros la fracción con numerador 5 y denominador 7»
49. Con join-word y use-dfrac: $\spmQ[join-word=enteros,use-dfrac=true]{2}{5}{7}$.

% doc: «1 entero 5 séptimos»
50. Con todas las llaves: $\spmQ[join-word=entero,use-spnZ=false,use-dfrac=true,wrap-parens=true,read-parens=false]{1}{5}{7}$.

% doc: «2 y la fracción con numerador 5 y denominador 7 más 1 y la fracción con numerador 1 y denominador 2»
51. Suma de números mixtos: $\spmQ{2}{5}{7} + \spmQ{1}{1}{2}$.

\section*{G. Comando sprazon}

% doc: «3 es a 4»
52. Dos enteros: $\sprazon{3}{4}$.

% doc: «a es a b»
53. Con letras: $\sprazon{a}{b}$.

% doc: «x más 1 es a 2»
54. Con una expresión: $\sprazon{x+1}{2}$.

% doc: «12 es a 18»
55. Con cifras grandes: $\sprazon{12}{18}$.

% doc: «3 es a 4»
56. Con silent false, el valor por defecto: $\sprazon[silent=false]{3}{4}$.

% doc: «»
57. Con silent true: $\sprazon[silent=true]{3}{4}$.

% doc: «3 es a 4»
58. Con use-dfrac true: $\sprazon[use-dfrac=true]{3}{4}$.

% doc: «3 es a 4»
59. Con inline true: $\sprazon[inline=true]{3}{4}$.

% doc: «a es a b»
60. Con inline true y letras: $\sprazon[inline=true]{a}{b}$.

% doc: «3 es a 4»
61. Con inline y use-dfrac: $\sprazon[inline=true,use-dfrac=true]{3}{4}$.

% doc: «»
62. Con silent e inline: $\sprazon[silent=true,inline=true]{3}{4}$.

% doc: «3 es a 4 es igual a 6 es a 8»
63. Dos razones iguales: $\sprazon{3}{4} = \sprazon{6}{8}$.

\section*{H. Comando spProp}

% doc: «3 es a 4 como 6 es a 8»
64. Cuatro enteros: $\spProp{3}{4}{6}{8}$.

% doc: «a es a b como c es a d»
65. Con letras: $\spProp{a}{b}{c}{d}$.

% doc: «x es a 2 como x más 1 es a 3»
66. Con expresiones: $\spProp{x}{2}{x+1}{3}$.

% doc: «3 es a 4 como 6 es a 8»
67. Con silent false, el valor por defecto: $\spProp[silent=false]{3}{4}{6}{8}$.

% doc: «»
68. Con silent true: $\spProp[silent=true]{3}{4}{6}{8}$.

% doc: «3 es a 4 como 6 es a 8»
69. Con use-dfrac true: $\spProp[use-dfrac=true]{3}{4}{6}{8}$.

% doc: «3 es a 4 como 6 es a 8»
70. Con inline true: $\spProp[inline=true]{3}{4}{6}{8}$.

% doc: «a es a b como c es a d»
71. Con inline true y letras: $\spProp[inline=true]{a}{b}{c}{d}$.

% doc: «3 es a 4 como 6 es a 8»
72. Con inline y use-dfrac: $\spProp[inline=true,use-dfrac=true]{3}{4}{6}{8}$.

% doc: «»
73. Con silent e inline: $\spProp[silent=true,inline=true]{3}{4}{6}{8}$.

\section*{I. En contexto}

% doc: «3 es a 4 flecha doble derecha 3 es a 4 como 6 es a 8»
74. Razón y proporción en una misma fórmula: $\sprazon{3}{4} \Rightarrow \spProp{3}{4}{6}{8}$.

% doc: «a es a b como c es a d»
75. En fórmula destacada:
\[ \spProp[use-dfrac=true]{a}{b}{c}{d} \].

\end{document}
